\documentclass[12pt]{article}
\usepackage[utf8]{inputenc}
\usepackage[T1]{fontenc}
\usepackage{hyperref}
\title{G\"uil\'a Zapotec Language Resources for K-12 Curriculum Development}
\author{Compiled by Claude Sonnet 5.5 (claude-sonnet-5-5) for Andr\'as Kornai}
\date{}
\begin{document}
\maketitle

\section{Introduction}
G\"uil\'a Zapotec (ISO 639-3 ztu, Glottocode guil1236), also called Diza, is spoken in San Pablo G\"uil\'a, Tlacolula District, Oaxaca, Mexico. Ethnologue gives about 9,500 speakers (1990 census), and Wikitongues rates it as vigorous. The searches found no G\"uil\'a-specific dictionary or school materials. The resources below are mostly descriptive or scripture-related.

\section{Primary Educational Resources}
None found. No primers, textbooks or school programmes were located.

\section{Linguistic and Reference Materials}
\begin{itemize}
\item Wikipedia overview, which cites a 1997 thesis on verb morphology by Ausencia L\'opez Cruz and an online dictionary of the closely related San Dionisio Ocotepec Zapotec (Broadwell and Martinez 2014): \url{https://en.wikipedia.org/wiki/G%C3%BCil%C3%A1_Zapotec}
\item AILLA: notes on the historical development of San Pablo G\"uil\'a Zapotec: \url{https://islandora-ailla.lib.utexas.edu/islandora/object/ailla%3A243794}
\item Omniglot page on G\"uil\'a Zapotec (Diza): \url{https://omniglot.com/writing/zapotec-guila.htm}
\item Numeral-system description (Max Planck channumerals database): \url{https://lingweb.eva.mpg.de/channumerals/Zapotec-Guila.htm}
\end{itemize}

\section{Religious and Community Resources}
\begin{itemize}
\item Scripture Earth listing: \url{https://www.scriptureearth.org/00eng.php?iso=ztu}
\item Wikitongues language page: \url{https://wikitongues.org/languages/ztu/}
\end{itemize}

\section{Digital and Audio Resources}
\begin{itemize}
\item Global Recordings Network audio: \url{https://globalrecordings.net/en/language/ztu}
\end{itemize}

\section{Curriculum Development Context}
The language is vigorous, so community speakers could write and review materials. The related San Dionisio Ocotepec dictionary may be a useful reference, but it is a different variety. A curriculum would need original materials developed with the community.
\end{document}